\section{Glosario de Términos}
Con el fin de asegurar la trazabilidad documental y la claridad metodológica para todos los involucrados, se definen los siguientes términos, siglas y acrónimos utilizados en este plan:
% Puedes usar una lista descriptiva aquí:
\begin{itemize}
\item
  \textbf{SQAP (Software Quality Assurance Plan / Plan de Aseguramiento
  de Calidad de Software):} Documento formal bajo el cual se rigen los
  procedimientos, métricas, roles y estándares destinados a certificar
  que el producto cumple con los niveles de aceptación pactados.
\item
  \textbf{SCM (Software Configuration Management / Gestión de la
  Configuración de Software):} Disciplina que consiste en aplicar
  procedimientos técnicos y organizativos para dirigir y vigilar la
  evolución, versiones, líneas base e integridad del código fuente y sus
  configuraciones.
\item
  \textbf{ORM (Object-Relational Mapping / Mapeo Objeto-Relacional):}
  Herramienta o técnica de software (como Eloquent en Laravel) que mapea
  la estructura relacional de la base de datos a un modelo orientado a
  objetos, evitando escribir consultas SQL manuales.
\item
  \textbf{HU (Historia de Usuario):} Unidad de medida funcional en
  metodologías ágiles que describe un requerimiento desde la perspectiva
  del usuario que interactuará con el sistema.
\item
  \textbf{XSS (Cross-Site Scripting / Secuencias de Comandos en Sitios
  Cruzados):} Vulnerabilidad de seguridad web en la cual un usuario
  inyecta scripts de código dañino en los campos de datos de entrada,
  afectando la navegación y seguridad de otros usuarios.
\item
  \textbf{CI/CD (Continuous Integration and Continuous Delivery /
  Integración y Despliegue Continuo):} Metodología de desarrollo en la
  cual los cambios de código se prueban de forma automática (CI) y se
  despliegan automáticamente a entornos de prueba o producción (CD).
\item
  \textbf{Soft-delete (Borrado Lógico):} Procedimiento en el cual los
  registros eliminados de la base de datos no se suprimen de los
  sectores físicos de almacenamiento, sino que se les asigna una marca
  temporal (campo \texttt{deleted\_at}), de forma que queden ocultos de
  las consultas normales pero se conserve el registro histórico para
  auditoría.
\item
  \textbf{SLA (Service Level Agreement / Acuerdo de Nivel de Servicio):}
  Compromiso formal que define el tiempo límite establecido para atender
  (dar respuesta inicial) y solucionar (resolver el defecto) un error de
  software según su nivel de prioridad o severidad.
\item
  \textbf{PR (Pull Request / Petición de Integración):} Mecanismo de
  colaboración y control en Git que avisa a los revisores de código que
  el trabajo de desarrollo de una rama de características está completo
  y listo para ser analizado e integrado a la rama principal.
\end{itemize}